%%%%%%%%%%%%%%%%%%%%%%%%%%%%%%%%%%%%%%%%%%%%%%%%%%%%%%%%%%%%%%%%
\section{Discussion of Results}
\label{sec:discussion_results}

The results answer the central research question in a bounded form.
They show that heterogeneous hybrid co-simulation can be implemented in one Python framework when the participating models are wrapped behind a common component interface and when the orchestration layer provides structural analysis, event handling, and state transfer.
The contribution is the integration architecture that connects different simulation tools and model fidelities in one executable workflow.

CoFMPy is the closest Python-native peer for this thesis~\cite{friedrich_cofmpy_2025}.
It provides \ac{FMI}-based co-simulation with Jacobi and Gauss--Seidel master algorithms, fixed-point algebraic-loop handling, and a Python \ac{FMU} proxy for non-\ac{FMU} components.
The overlap with \syssimx{} lies in Python-based \ac{FMU} orchestration and algebraic-loop handling.
The contribution boundary of \syssimx{} lies in the shared component interface for \ac{FMU}, OpenSim, and FEM backends, hybrid event localization with state restoration, runtime model switching, and Newton-type interface-Jacobian loop solving.
These mechanisms are verified in Chapter~\ref{chap:implementation} and combined in Chapter~\ref{chap:case_study}.

\paragraph{RQ1. Heterogeneous Integration.}
The first research question asked how heterogeneous models from equation-based, musculoskeletal, and finite-element tools can be represented through a unified interface.
The implementation answers this question through the port system and the \texttt{CoSimComponent} abstraction.
\acp{FMU}, OpenSim models, \ac{FEM} models, and multi-model components are exposed through the same input, output, parameter, state, and stepping concepts.
Backend differences remain inside wrapper-specific code.
The system-level orchestration can therefore operate on a common interface.
This result is consistent with the role of \ac{FMI} as an interoperability standard~\cite{modelica_association_project_fmi_functional_2014,modelica_association_project_fmi_functional_2024}.
It extends the workflow beyond \ac{FMU}-only coupling by including tool-specific OpenSim and \ac{FEM} components.

\paragraph{RQ2. Dependency-Aware Orchestration.}
The second research question concerned dependency analysis and execution ordering.
The structural-analysis implementation derives persistent metadata from the registered components, signal connections, and direct-feedthrough information.
This metadata is then reused by coupled initialization, master-algorithm execution, and algebraic-loop handling.
The important point is that the execution order is not encoded manually in each scenario.
It is derived from the assembled system structure.
The feature-level verification in Chapter~\ref{chap:implementation} therefore supports the claim that \syssimx{} can identify instantaneous dependencies and algebraic-loop blocks before the simulation is run~\cite{arnold_error_2013}.
The need to detect and resolve such cyclic dependencies follows from the classical analysis of modular simulator coupling, which shows that algebraic loops between  direct-feedthrough subsystems destabilize modular integration unless the interface variables are solved iteratively at the communication points~\cite{kubler_two_2000}.
\syssimx{} follows this iterative-coupling tradition and resolves detected loops with the interface-Jacobian solver rather than by altering the coupled dynamics.
The implementation operates on component-level dependency information rather than full symbolic equations.
This scope is sufficient for the metadata available in the framework.

\paragraph{RQ3. Hybrid Event Handling.}
The third research question addressed hybrid phenomena such as contact and discrete state changes.
The hybrid algorithm answers this question by combining trial stepping, rollback, event localization, dense-time event handling, and event dispatch.
The verification scenarios in Chapter~\ref{chap:implementation} show that this mechanism can localize zero-crossing events, process cascaded events at the same physical time, and check simultaneous event handling for commutativity.
The controlled-pendulum contact scenario then shows the same mechanisms in the case-study context.
The contact event is localized, routed to the controller, and used to reset the PID integrator.
This connects the implementation feature with the coupled system workflow.
This result follows known requirements for hybrid co-simulation.
Rollback, event ordering, and time representation are central issues in this setting~\cite{broman_requirements_2015,cremona_hybrid_2019}.

\paragraph{RQ4. Runtime Model Switching.}
The fourth research question asked how alternative subsystem models can be switched at runtime.
The multi-model component answers this through a fixed external interface and an internal active-model state.
The \texttt{MasterPendulum} uses this concept to switch between FMU, OpenSim, and FEM pendulum models while the surrounding closed-loop system sees one pendulum component.
The switching verification shows that the exposed rigid-body quantities remain continuous at the switch points.
The performance comparison also shows the consequence of this design choice.
Switching into the FEM model reconstructs the finite-element state from reduced rigid-body variables.
It therefore cannot reproduce deformation history that was not simulated before the switch.
The interface remains continuous.
The projection still loses internal FEM history.
This is the modeling cost of switching between representations with different internal state spaces.

\paragraph{RQ5. Case-Study Behavior.}
The fifth research question concerned representative system behavior and benchmark scenarios.
The controlled-pendulum case study provides the system-level evidence.
The baseline scenario verifies the FMU-only closed-loop co-simulation against a monolithic OpenModelica reference.
The decreasing angle error for smaller macro step sizes confirms the expected macro-step dependence.
This behavior is consistent with the error mechanisms of explicit co-simulation~\cite{gomes_co-simulation_2019,arnold_error_2013}.
The contact scenario verifies the combined workflow against a rigid-contact reference and shows that the framework can combine heterogeneous plant models, hybrid contact events, PID reset, and runtime switching in one simulation.
It uses the rigid-contact reference as a numerical comparator and demonstrates the workflow on a controlled hybrid system.

The performance benchmark gives the practical interpretation of runtime model switching.
For the reported configuration, restricting \ac{FEM} execution to the contact-relevant phase produced an end-to-end speedup of about \(1.6\times\).
The \ac{FEM} solver itself benefited more strongly than the total run time.
This difference shows that the switched configuration reduces expensive model evaluations, while orchestration and non-\ac{FEM} components still contribute to the total cost.
The benchmark therefore supports runtime model switching for the selected scenario.
It does not imply a general speedup for all multi-fidelity simulations.

\paragraph{Relation to Multi-Fidelity Simulation.}
The runtime model switching of \syssimx{} is a hierarchical fidelity-management strategy rather than a surrogate-based one, because it switches between executable subsystem models instead of fusing their outputs into a single approximation~\cite{fernandez-godino_review_2023}.
In the model-management taxonomy of Peherstorfer et al., the contact-driven activation of the \ac{FEM} model is closest to a filtering strategy, in which an inexpensive model determines when the expensive model is required~\cite{peherstorfer_survey_2018}.
The closest operational precedent is the interest-region switching of Choi et al., where a selection model activates a high-fidelity submodel only inside a defined region and transfers the state to the newly activated model~\cite{Choi2014}.
The reported end-to-end speedup of \(1.61\times\) is comparable to the \(1.21\times\)--\(1.25\times\) end-to-end speedups reported by Choi et al., and the benefit is similarly bounded by orchestration overhead and switching frequency.
These methods target outer-loop tasks such as optimization and uncertainty propagation, whereas \syssimx{} applies fidelity switching inside a time-domain hybrid co-simulation loop with heterogeneous tool backends.

Taken together, the results support the thesis contribution.
The framework mechanisms were verified in isolation and then combined in the controlled-pendulum case study.
The next section defines the scope in which this claim holds.
